\chapter*{The Artificial-Intelligence Map for This Course}
\addcontentsline{toc}{chapter}{The Artificial-Intelligence Map for This Course}
\markboth{The Artificial-Intelligence Map}{}

The laboratory material spans several AI techniques that solve different kinds of problems. A useful first question is not ``Which algorithm is best?'' but ``What kind of information does the problem give me, and what must the system produce?''

\begin{mltable}
\begin{tabularx}{\linewidth}{@{}>{\bfseries\leavevmode\color{MLNavy}}p{0.19\linewidth}p{0.24\linewidth}X@{}}
\toprule
\rowcolor{MLPaleBlue!92}
Problem signal & Main techniques in the labs & Typical output \\
\midrule
Connectivity only & BFS, DFS & A reachable path or traversal order \\
Edge costs & Uniform-Cost Search & A least-cost path \\
Heuristic estimate & Greedy Best-First Search & A goal-directed path \\
Cost + heuristic & A* Search & A path balancing known cost and estimated remaining cost \\
Imprecise linguistic rules & Fuzzy Logic, FIS & A crisp control value from fuzzy rules \\
Input-output examples & ANFIS & A learned Sugeno-type fuzzy model \\
Robot state and feedback & Kinematics, PID, A* planning & Motion, path, or control response \\
Text or images containing text & Embeddings, OCR, simple NLP rules & Similarity, labels, summaries, translations, keywords \\
Candidate solution population & Genetic Algorithm & An optimized or near-optimized solution \\
\bottomrule
\end{tabularx}
\end{mltable}

\section*{Search methods at a glance}

\begin{mlfigure}
\begin{tikzpicture}[
  q/.style={draw=MLNavy,fill=MLNavy,text=white,rounded corners=2mm,text width=13.4cm,minimum height=0.85cm,align=center,font=\small\bfseries},
  b/.style={draw=MLBlue,fill=MLPaleBlue,rounded corners=2mm,text width=2.7cm,minimum height=1.35cm,align=center,font=\scriptsize\color{MLNavy}},
  arr/.style={-{Stealth[length=2mm]},thick,draw=MLTeal}
]
\node[q] (q) at (0,2.2) {What information is used to choose the next frontier node?};
\node[b] (bfs) at (-5.2,0) {\textbf{BFS}\\shallowest first\\FIFO queue};
\node[b] (dfs) at (-1.75,0) {\textbf{DFS}\\deepest first\\stack / recursion};
\node[b] (ucs) at (1.75,0) {\textbf{UCS}\\lowest $g(n)$\\priority queue};
\node[b] (gbfs) at (5.2,0) {\textbf{GBFS}\\lowest $h(n)$\\priority queue};
\node[b,below=1.0cm of ucs] (astar) {\textbf{A*}\\lowest $f(n)=g(n)+h(n)$};
\draw[arr] (q.south) -- (0,1.35);
\draw[draw=MLTeal,thick] (-5.2,1.35) -- (5.2,1.35);
\foreach \x in {bfs,dfs,ucs,gbfs} {\draw[arr] (\x.north) -- ++(0,0.85);}
\draw[arr] (ucs.south) -- (astar.north);
\end{tikzpicture}
\end{mlfigure}

\begin{keyidea}
BFS and DFS differ mainly in frontier order. UCS adds path cost. Greedy Best-First Search adds a heuristic but ignores accumulated cost. A* combines both quantities.
\end{keyidea}
